Flash floods are the most destructive natural hazard in Himachal Pradesh (HP), India,
causing over 400 fatalities and \$1.2 billion in losses in the 2023 monsoon season alone.
Existing susceptibility maps for the region are limited to a single river basin,
use random train-test splits that inflate accuracy estimates,
and provide only point-estimate risk scores with no quantified uncertainty.
This study presents the first state-wide, multi-basin flash flood susceptibility assessment
for HP that addresses all three shortcomings simultaneously.

We constructed a six-year flood inventory (2018--2023) from Sentinel-1 SAR imagery
processed in Google Earth Engine using a seasonal minimum backscatter approach,
supplemented by terrain-filtered flood sampling (slope\,$<$\,15°, within 2\,km of channels)
yielding 3,000 georeferenced flood occurrences across HP.
Twelve conditioning factors spanning terrain morphology, hydrology, land cover, soil,
and rainfall were assembled at 30\,m resolution after multicollinearity screening
(one pair removed: $|r|=0.97$; remaining factors with max VIF\,=\,1.84).
Four model families were evaluated: Random Forest (RF), XGBoost, LightGBM,
and a stacking ensemble, serving as pixel-based baselines.
As the primary methodological contribution, we trained a Graph Neural Network (GraphSAGE)
on a directed watershed connectivity graph (460 sub-watersheds, 1,700 edges),
explicitly encoding upstream--downstream flood propagation that pixel-based models cannot capture.
All models were validated with leave-one-basin-out spatial block cross-validation
to avoid the 5--15\% AUC inflation typical of random splits.
As a second novel contribution, we applied split-conformal prediction (MAPIE)
to produce the first susceptibility maps with statistically guaranteed 90\% coverage intervals,
enabling prioritisation of high-risk, high-certainty zones for emergency management.

The GNN achieved a mean spatial-CV AUC of 0.995\,$\pm$\,0.004, outperforming the best
pixel-based model (stacking ensemble: AUC\,=\,0.901\,$\pm$\,0.004) and surpassing the
benchmark reported by \citet{saha2023} (AUC\,=\,0.88) for the same study region.
Conformal coverage analysis showed that the 90\% prediction intervals achieved
82.9\% empirical coverage on the held-out temporal test set, with systematic
undercoverage in high-risk areas (45--59\%) attributable to SAR label noise.
SHAP analysis identified elevation, plan curvature, and slope as the three most globally
important susceptibility drivers, with marked spatial variation across districts.
High-susceptibility zones overlay approximately 1,457\,km of national and state highways,
2,759 bridges, and 4 major hydroelectric installations,
representing critical infrastructure exposure
that is directly actionable by HP State Disaster Management Authority.

\noindent\textbf{Plain language summary:}
We created the most complete flood risk map yet produced for Himachal Pradesh
by teaching a machine-learning model that water flows downhill through connected river valleys
(not just through isolated grid cells), and by reporting not only \textit{where} flooding is likely
but \textit{how confident} the model is in each prediction.
